\scanpage{54}
un Pb. est celui où $k(\handwritten{\eta})=L$ est de car.$p$ (cas d'égales caractéristiques).
On peut alors penser à attaquer le Pb via la théorie de Dieudonné (supposan\deleted{\unsure{s}}t résolu le Pb 1) et via le

\uline{Problème} n\textsuperscript{\uline{o}} ... Soit $S$ un schéma noethérien normal \deleted{de car} intègre d\deleted{\unsure{a}}e car.
$p$, de \deleted{\unsure{o}}point générique $\handwritten{\eta}$. Est-il vrai que le foncteur de restriction
\[
\text{Cris loc lib }(S)\handwritten{\longrightarrow}\text{ Cris loc lib }(\handwritten{\eta})
\]
(ou du moins le foncteur analogue pour F-cristaux) est pleinement fidèle ?

On est ramené facilement au cas où $S$ est le spectre d'un anneau de
\deleted{\unsure{g}}la forme $k\handwritten{[[}t\handwritten{]]}$, avec $k$ alg. clos, auquel cas l'énoncé s'explicite en
termes de Modules avec connexions formelles .... Il semble que si la réponse au Pb précédent est affirmative, \deleted{ainsi que} et si on sait faire la
théorie de Dieudonné\unsure{n} sur un anneau \uline{parfait} (avec \uline{équivalence} de catégories
\deleted{\unsure{q}}alors on aura démonstration du th. de Tate en égales car.
% archive sheet number: 53

